\chapter{四捨五入と和の順序と丸め誤差}
    \begin{shadebox}
        $N,B\in\naturalNumbers,x_i\in\integers\;(i=0,1,\dots,N)$ とする。
        丸め誤差は一様分布するものとする。
        このとき $x_i$ の総和の下位 $B$ ビットを四捨五入した値 $S_\mathrm{sr}$ と $x_i$ の下位 $B$ ビットを四捨五入した後の総和 $S_\mathrm{rs}$ の誤差の期待値はともに 0 である。
        標準偏差については，後者が前者の $\sqrt{N}$ 倍である。
    \end{shadebox}
    ここでは 2 進表現を前提とするが，他の進数表現でも同じ結果が得られる。
    \begin{proof}
        \quad\par
        $S_\mathrm{sr}$ の丸め誤差は $\{-2^{B-1}+1,\dots,2^{B-1}\}$ に一様分布する。
        $x_i$ の下位 $B$ ビットを四捨五入した値の丸め誤差は $\{-2^{B-1}+1,\dots,2^{B-1}\}$ に一様分布する。
        これの和が $S_\mathrm{rs}$ の丸め誤差であるから，分散が $N$ 倍，標準偏差が $\sqrt{N}$ 倍になる。
    \end{proof}